\chapter{The Model Zoo I: Trees and MLPs}
\label{ch:treesmlp}

Parts I--IV built the world the models live in; the next four chapters walk
the models themselves. We take each family in turn: what it computes, the
inductive bias that decides where it wins, its measured record in this
campaign, and the operational details that textbooks omit. We begin with
the two families that ignore time entirely --- and that between them
powered the winning solutions of several competitions in this genre.

\section{Gradient-boosted trees}

\subsection{What it computes}

XGBoost/LightGBM fit an additive ensemble
$\hat y(x) = \sum_{m=1}^{M} f_m(x)$ of shallow regression trees, each fit
to the gradient of the loss at the current ensemble's predictions. Each
tree partitions feature space with axis-aligned splits; depth $d$ trees can
express interactions of order at most $d$.

\subsection{Why it is always the first model}

\begin{itemize}
\item \textbf{Nothing to tune before signal appears.} Reasonable defaults
  produce a competent model in one shot; neural nets need scaling,
  initialization, and schedule choices before they beat zero.
\item \textbf{Invariances that match tabular data:} monotone
  transformations of features do not matter; scales do not matter; nulls
  are handled natively; heavy tails only shift split points.
\item \textbf{It doubles as a measurement instrument:} SHAP importances
  drive the selection loop of Chapter~\ref{ch:selection}, and a fixed XGB
  harness is our standard ablation meter for feature blocks
  (Chapter~\ref{ch:featurepatterns}) --- fast, deterministic given a seed,
  and sensitive enough to resolve $10^{-4}$-scale effects with fixed data.
\end{itemize}

Its measured record here: $+0.0059$ on 200 recent days, $+0.0068$ on 280
--- competitive with far more elaborate architectures --- and a permanent
seat in the final ensemble because tree errors decorrelate beautifully
from neural errors (Chapter~\ref{ch:ensembling}).

\subsection{The blind spots, measured}

\begin{itemize}
\item \textbf{No temporal state.} A tree sees a row. All sequence context
  must arrive as features --- which is why Chapter
  \ref{ch:featurepatterns}'s engineered blocks exist. Where recurrent
  models gained $+0.003$--$0.005$ from online refit, trees as we ran them
  are static (retraining nightly is possible but a different, heavier
  operation).
\item \textbf{Smooth interactions cost depth.} The symbolic-combo result
  is a direct measurement: features like $\min(f_a, f_b)$, each parent
  already available, added $+0.0005$ --- structure a depth-6 forest
  demonstrably had not synthesized from 34 million rows. Axis-aligned
  splits approximate diagonal or multiplicative structure expensively;
  at low SNR the boosting process spends its statistical budget elsewhere
  first.
\item \textbf{Extrapolation is flat.} Outside the training hull,
  predictions saturate at leaf values. Under drift this is half blessing
  (no wild extrapolations), half curse (no adaptation).
\end{itemize}

\subsection{Operational notes}

Configuration that mattered here: \code{tree\_method=hist} with
\code{max\_bin=128} (accuracy indistinguishable, memory and speed much
better at this scale); \code{min\_child\_weight} $\sim$10 and subsampling
(row 0.8 / column 0.5) as the main overfitting brakes; early stopping on
the \emph{weighted} eval metric with weights passed to the eval set
(Chapter~\ref{ch:metric}); $n \sim 1500$ rounds offered but early stopping
typically firing near 100--170. One platform trap that cost us an
afternoon: mixing xgboost and torch in one process on macOS doubles up the
OpenMP runtime --- deadlocks at high thread counts, segfaults on pickle;
we run XGB single-threaded inside evaluation harnesses and never pickle
boosters alongside torch (prediction dumps as \code{.npz} instead).

\section{The MLP}

\subsection{What it computes, and the DRW lesson}

Dense layers over the feature vector; no recurrence, no attention, no
notion of time. It is tempting to dismiss --- our own first result was a
plain MLP on all 130 columns scoring $\bm{-0.0026}$, worse than predicting
zero. The DRW crypto winner's \emph{protagonist} model was a 3-layer MLP
scoring at the top of a 1{,}000-team leaderboard. Same family; the
difference was everything upstream: he fed it $\sim$40 forensically
curated features plus symbolic combos plus AE factors, while our early MLP
drank from the raw firehose.

\begin{keybox}[: the MLP is a measurement of your features]
An MLP has just enough capacity to exploit clean features and no inductive
machinery to rescue dirty ones. If your MLP is far below your XGB, your
feature set relies on tree-style robustness (nulls, outliers, monotone
weirdness) --- clean it. If your MLP approaches or beats trees, your
features are doing the work and \emph{every} downstream architecture will
benefit. Track the XGB--MLP gap as a feature-quality dashboard.
\end{keybox}

\subsection{Training an MLP at low SNR: the DRW recipe}

Three of his choices generalize and are cheap to adopt:

\begin{itemize}
\item \textbf{Loss = $0.6\,\mathrm{MSE} + 0.4\,(1 - \rho_{\text{Pearson}})$.}
  Pure MSE gradients at $R^2 = 0.01$ are dominated by irreducible noise;
  the correlation term rewards co-movement --- the controllable part ---
  and empirically stabilized his training. (Where the metric \emph{is}
  weighted $R^2$ and batches are whole days, training on $1-\Rw$ directly
  is the equivalent move; Chapter~\ref{ch:metric}.)
\item \textbf{SGD over Adam.} His MLPs generalized better under plain SGD;
  a commenter reported Muon better still. Adaptive optimizers' per-parameter
  step sizes are quick to fit noise structure; at low SNR, dumber and
  slower optimization is a regularizer. Our recurrent models kept AdamW
  \emph{for the offline phase} but the finding transfers as: treat the
  optimizer as a regularization hyperparameter, not a solved default.
\item \textbf{Robustness over cleverness.} His stated post-mortem: added
  regularization could not recover a score that feature drift had taken;
  re-auditing features (the recycling step, Chapter~\ref{ch:selection})
  could. When an MLP degrades on newer data, suspect the features before
  the weights.
\end{itemize}

\subsection{Where the MLP sits in a panel pipeline}

Even without temporal machinery, the MLP has two structural roles:

\begin{enumerate}
\item \textbf{Ensemble diversifier}: its errors decorrelate from both
  trees (different function class) and RNNs (no temporal state), making it
  a candidate stream even at a lower solo score
  (Chapter~\ref{ch:ensembling}'s admission rule).
\item \textbf{Head for engineered temporal features}: an MLP over
  signature features, atlas-selected lags, and rolling statistics is a
  legitimate sequence model in disguise --- all the temporal work moved
  into the features. Our \code{mlp\_sig} variant (MLP over path-signature
  features, Chapter~\ref{ch:signatures}) scored $+0.0066$ online ---
  above XGB, below the RNNs --- with the entire temporal burden carried
  by its inputs.
\end{enumerate}

\section{Trees versus MLP versus both}

The genre's empirical regularity, visible across this competition, DRW,
and the public solutions we studied: \textbf{tree ensembles and neural
models finish close together, and their average finishes ahead of either.}
The reason is not mystical: they misfit in different directions ---
axis-aligned versus smooth, interpolating versus extrapolating,
feature-robust versus feature-hungry. You are not choosing a champion; you
are hiring complementary employees. The only real decision is which gets
your tuning hours, and the answer follows from the feature-quality
dashboard above: dirty features $\to$ trees first; clean features $\to$
neural first; either way, both ship.

\begin{drillbox}
Reproduce the feature-quality dashboard on your own data: fix one feature
matrix, train (a) default XGBoost, (b) a 3-layer MLP with standardized
inputs, MSE + Pearson loss, SGD. Record the gap. Then add your five best
engineered features (Chapter~\ref{ch:featurepatterns}'s drill) and record
the gap again. The XGB score should move a little; the MLP score should
move a lot --- if it doesn't, your engineered features encode robustness,
not information.
\end{drillbox}
